% chktex-file 9
% chktex-file 17
% chktex-file 36

\subsection*{1. 平均自由徑}
\subsection*{2. 自由路徑分布}
~\cite{https://doi.org/10.1002/andp.18521620505}的論文說明粒子通過均勻介質時，存活下來的機率成指數的衰減，因此假設粒子前進 $x$ 後存活下的機率符合方程：
\begin{equation}
    \frac{d P \left(x\right)}{dx}=-\frac{1}{\lambda} P(x)
\end{equation}
解出此方程式的解為
\begin{equation}
    P(x) = e^{\frac{-x}{\lambda}}
\end{equation}
藉此可求出粒子剛好在距離 \(x\) 到 \(x + dx\) 之間發生第一次碰撞的機率：
\begin{equation}
    f(x)=P(x)\cdot \frac{1}{\lambda }=\frac{1}{\lambda }e^{-x/\lambda }\quad (x\ge 0)
\end{equation}

轉成累積分佈函數(CDF)後
\begin{equation}
    F(s)=\int _{0}^{s}\frac{1}{\lambda }e^{-x/\lambda }\,dx=1-e^{-s/\lambda }
\end{equation}
顯然 $F(s)$ 的值域嚴格落在 $[0,1)$。\\
求解 $F^{-1}(\zeta),\ \zeta \in [0,1)$：
\begin{align}
    \zeta &= 1-e^{-s/\lambda } \\
    s &= -\lambda \ln\left(1 - \zeta\right)
\end{align}

因此若 $R \in (0,1]$ 唯一隨機隨機數，則可透過
\begin{equation}
    s = F^{-1}\left(R\right) = -\lambda \ln \left(R\right)
\end{equation}
得出粒子的自由路徑分布。